%----------------------------------------------------------%
%          TEMPLATE - REVISTA SEMINA (26/01/2018)          %
%----------------------------------------------------------%

% CLASSE DE DOCUMENTO
\PassOptionsToPackage{bookmarks=false}{hyperref}
\documentclass[article,10pt,twoside,a4paper,twocolumn,hidelinks,english,brazil]{abntex2}

% PACOTES
\usepackage[utf8]{inputenc}
\usepackage[T1]{fontenc}
\usepackage{helvet}
\usepackage{indentfirst}
\usepackage{color}
\usepackage{fancyhdr,parskip}
\usepackage[margin=1in]{geometry}
\usepackage{graphicx}
\usepackage{microtype}
\usepackage{soulutf8}
\usepackage{mathptmx}
\usepackage{titling}
\usepackage{etoolbox}
\usepackage{titlesec}
\usepackage{ragged2e}
\usepackage{float}
\usepackage[labelfont=bf,format=plain,justification=justified,singlelinecheck=false]{caption}
\usepackage[bookmarks=false]{hyperref}
\usepackage{helvet}
\usepackage[alf,bibjustif,abnt-etal-list=0,abnt-etal-cite]{abntex2cite}
%\usepackage[brazilian,hyperpageref]{backref}
%\usepackage[alf,bibjustif,abnt-etal-list=0,abnt-etal-cite]{abntcite}

\usepackage{amsmath,amssymb,amsthm}
\usepackage{graphicx}
\usepackage{lastpage}
\usepackage{multirow}


%\addbibresource

%\usepackage[alf]{abntex2cite}
%\usepackage{amsmath}
\usepackage{wrapfig}
%\usepackage[hang]{footmisc}
%%%%incluido
%\usepackage[num,overcite]{abntex2cite}
%\citebrackets[]
\usepackage{url}  % pacote para quebrar linha da URL

% AGREGA LA IMAGEN EN LA PARTE SUPERIOR IZQUIERDA
\usepackage{eso-pic}
\AddToShipoutPictureBG*{%
  \AtPageUpperLeft{%
    \raisebox{-\height-\baselineskip}{% Mueve la imagen 0.5 cm hacia abajo
      \hspace*{0.4cm} % Mueve la imagen 1 cm a la izquierda
      \includegraphics[width=2cm]{C:/Users/ADMON/Documents/Universidad/logo.png}%
    }%
  }
}


% MARGENS ESQUERDA, DIREITA, SUPERIOR E INFERIOR
\setlrmarginsandblock{2.5cm}{2cm}{*}
\setulmarginsandblock{2cm}{2.3cm}{*}
\checkandfixthelayout

% OUTROS ESPACAMENTOS
\setlength{\droptitle}{-1.3cm}
\setlength{\columnsep}{0.5cm}
\setlength{\parindent}{0.5cm}
\setlength{\parskip}{0.0cm}

% LINHA HORIZONTAL
\def\Vhrulefill{\leavevmode\leaders\hrule height 0.7ex depth \dimexpr0.5pt-0.7ex\hfill\kern0pt}

% NOTAS DE RODAPE
\thanksmarkseries{arabic}
\setlength\thanksmarkwidth{.5em} %adicionei
\renewcommand{\footnoterule}{
    \kern -3pt
    \hrule width 4.2cm height 0.5pt
    \kern 2.6pt
}
\let\oldmakethanksmark\makethanksmark
\renewcommand{\makethanksmark}{\oldmakethanksmark\small}
\let\oldthanks\thanks
\renewcommand{\thanks}[1]{\oldthanks{\small #1}}


% SECOES E SUBSECOES
\titleformat*{\section}{\fontsize{11pt}{13.2pt}\bfseries}
\titleformat*{\subsection}{\itshape}{\fontsize{11pt}{1.2pt}}
\titleformat*{\subsubsection}{\itshape}{\fontsize{11pt}{1.0pt}}

% REFERENCIAS
\apptocmd{\thebibliography}{\justifying}{}{}

% ----ALGUNS DESTES DADOS SERÃO ADEQUADOS PELA EDITORAÇÃO CIENTÍFICA  -----------%

% INFORMACOES DO ARTIGO E DA REVISTA

\newcommand{\paginainicial}{1} % página inicial




%----------autores tem que informar os dados AQUI----------
\newcommand{\autorcabecalho}{Author1, F.~A.; Author2, Z.~Z.; Author3, R.;  Author4, S.~W.} % nomes dos autores no cabecalho
\newcommand{\titulocabecalho}{} % titulo no cabecalho




%---------------------------------------------------------%

% PAGINA INICIAL
\setcounter{page}{\paginainicial}

% CABECALHO E RODAPE DA PRIMEIRA PAGINA
\makepagestyle{firstpage}
\makeoddhead{firstpage}{\hfill\fontfamily{phv}\scriptsize 35$^{\circ}$ Simposio Internacional de Estadística 2026 \\[0.15cm]
\hfill Fecha: Julio 14 - Julio 17}{}{}



\makeoddfoot{firstpage}{}{\fontfamily{phv}
\footnotesize \selectfont \hspace{1cm} \Vhrulefill \hspace{0.65cm} \thepage \\ 35$^{\circ}$ Simposio Internacional de Estadística 2025 }







% CABECALHO E RODAPE DAS DEMAIS PAGINAS
\makepagestyle{text}
\makeevenhead{text}{}{\fontfamily{phv} \footnotesize \selectfont \autorcabecalho}{}
\makeoddhead{text}{}{\fontfamily{phv} \footnotesize \selectfont \titulocabecalho}{}
\makeevenfoot{text}{}{\fontfamily{phv} \footnotesize \selectfont \thepage \hspace{0.65cm} \Vhrulefill \hspace{1cm}~\\ Semina: Ciências Exatas e Tecnológicas, Londrina, v. \volume, n. \numero, p. \paginainicial-\pageref{LastPage}, \mes \ano}{}
\makeoddfoot{text}{}{\fontfamily{phv} \footnotesize \selectfont \hspace{1cm} \Vhrulefill \hspace{0.65cm} \thepage \\ Semina: Ciências Exatas e Tecnológicas, Londrina, v. \volume, n. \numero, p. \paginainicial-\pageref{LastPage}, \mes \ano}{}


%----------------------- OS AUTORES DEVEM INCLUIR AS INFORMAÇÕES ABAIXO-------------%

% TITULO % Gustavo 19 maio

\titulo{\fontsize{16pt}{19.2pt} \selectfont
    \textbf{Modelando la Curva de Rendimientos en Colombia: DNS y DNSS bajo la Lupa Bayesiana y Frecuentista}  % titulo
}
\tituloestrangeiro{\fontsize{12pt}{18pt} \selectfont
    \textit{\textbf{Modeling the Yield Curve in Colombia: DNS and DNSS Under Bayesian and Frequentist Lenses}} % titulo em inglês
}


% AUTORES: Nome completo. Sobre os demais dados use abreveaturas, se necessário, para manter todas as informações/descrições,  de cada autor, em apenas uma  linha.  




\autor{ \fontsize{13pt}{15.6pt} \selectfont
    Juan David Garro\thanks{Est. Estadistica , UNAL - Med, E-mail: kgonzalezd@unal.edu.co}; % primer autor
    Kevin Hair Gonzalez Duarte\thanks{Est. Ingeniería Administrativa , UNAL - Med, E-mail: jgarro@unal.edu.co}; % segundo autor}
    Monica Andrea Arango Arango\thanks{Profesora Titular, UNAL - Med, E-mail: marango@unal.edu.co}; % tercer autor
    \\
    %\colorbox[rgb]{1,1,0}{autores enviar  informações completas - Mini currículo em português}
   % Author 5\thanks{ Master student, Dept. Chemistry, USP, SP, Brazil; E-mail: author5@gmail.com}
}

%---------------------------------------------------------%

% DATA
\data{}

\begin{document}

\captionsetup{format=plain,justification=justified,singlelinecheck=false}

% DEFINICOES DO ARTIGO E PRE-TEXTUAIS
\selectlanguage{english}
\frenchspacing
% \nocite{*} % comentar esta linha se houver \cite

\pretextual
\SingleSpacing
\onecolumn

% TITULO E AUTORES
\maketitle
\vspace{-1.2cm}

%Inserto imagen a la página
%\thispagestyle{fancy}

% ESTILO DA PRIMEIRA PAGINA
\thispagestyle{firstpage}

%----------------------- OS AUTORES DEVEM INCLUIR AS INFORMAÇÕES ABAIXO-------------%
\hyphenation{} %  palavras sem quebras

% RESUMO E PALAVRAS-CHAVE EM INGLES % Gustavo 19 maio
\renewcommand{\resumoname}{\texorpdfstring{\fontsize{14pt}{16.8pt} \fontfamily{ptm} \selectfont \textbf{Resumen} \\[-0.3cm] \rule{0.91\textwidth}{0.5pt} \vspace{-0.5cm}}{Abstract}}
\begin{resumoumacoluna}
    \normalsize

Los Títulos de Tesorería (TES) constituyen el activo libre de riesgo y la columna vertebral del mercado de
capitales colombiano. Aunque los modelos de la familia Nelson-Siegel son ampliamente utilizados para modelar
su curva de rendimientos, la literatura tradicional suele fijar el parámetro de decaimiento ($\tau$) como una
constante, omitiendo la naturaleza dinámica de la estructura temporal de tasas. Para cubrir esta brecha empírica, 
este estudio contrasta rigurosamente la estimación de la curva cero-cupón mediante modelos de espacio de estados, 
evaluando metodologías frecuentistas frente a estimaciones bayesianas puras. Metodológicamente, se construye un 
panel histórico de tasas spot mediante el método de bootstrapping de Fama-Bliss en seis nodos de vencimiento
(entre 2 y 15 años). Sobre este panel, se extraen los factores latentes aplicando los modelos Dinámico de Nelson-Siegel
(DNS) y Nelson-Siegel-Svensson (DNSS) bajo dos paradigmas: (1) un enfoque frecuentista basado en el Filtro de Kalman, 
optimizando los parámetros $\tau$ a través de una búsqueda de grilla en ventana rodante, y (2) un enfoque Bayesiano 
mediante el algoritmo de Muestreo de Gibbs, que estima conjuntamente los factores $\beta$ y $\tau$. El desempeño de 
las metodologías se evalúa a través de métricas de bondad de ajuste. La hipótesis 
central sugiere que la estimación dinámica de $\tau$ reduce significativamente los errores de pronóstico. Se anticipa 
que el enfoque Bayesiano ofrezca una ventaja predictiva al capturar mejor la incertidumbre durante periodos volátiles,
mientras que el Filtro de Kalman destaca por su eficiencia computacional y ajuste competitivo. Esta investigación proyecta consolidar un marco analítico superior frente a los modelos estáticos tradicionales, proveyendo herramientas robustas para la valoración de activos, gestión de riesgos y toma de decisiones en Colombia.

\vspace{0.4 cm}
    \noindent \textbf{Palabras clave:} xxxxxxxxxxxxxxxx

\end{resumoumacoluna}


% RESUMO E PALAVRAS-CHAVE EM PORTUGUES % Gustavo 19 maio
\renewcommand{\resumoname}{\texorpdfstring{\fontsize{14pt}{16.8pt} \fontfamily{ptm} \selectfont \textbf{Abstract} \\[-0.3cm] \rule{0.91\textwidth}{0.5pt} \vspace{-0.5cm}}{Resumo}}
\begin{resumoumacoluna}
    \normalsize

The Treasury Bonds (TES) are the risk-free asset and the backbone of the Colombian capital market. 
Although the Nelson-Siegel family of models is widely used to model their yield curve, traditional 
literature often fixes the decay parameter ($\tau$) as a constant, overlooking the dynamic nature of
the term structure of rates. To fill this empirical gap, this study rigorously contrasts the estimation
of the zero-coupon curve through state-space models, evaluating frequentist methodologies against pure
Bayesian estimations. Methodologically, a historical panel of spot rates is constructed using the 
Fama-Bliss bootstrapping method at six maturity nodes (between 2 and 15 years). On this panel, latent
factors are extracted by applying Dynamic Nelson-Siegel (DNS) and Nelson-Siegel-Svensson (DNSS) models 
under two paradigms: (1) a frequentist approach based on the Kalman Filter, optimizing $\tau$ parameters 
through a grid search in a rolling window, and (2) a Bayesian approach using the Gibbs Sampling algorithm, 
which jointly estimates $\beta$ and $\tau$ factors. The performance of the methodologies is evaluated through 
goodness-of-fit metrics. The central hypothesis suggests that dynamic estimation of 
$\tau$ significantly reduces forecasting errors. It is anticipated that the Bayesian approach will offer a 
predictive advantage by better capturing uncertainty during volatile periods, while the Kalman Filter stands 
out for its computational efficiency and competitive fit. This research projects to consolidate a superior 
analytical framework compared to traditional static models, providing robust tools for asset valuation, risk 
management, and decision-making in Colombia.   

\vspace{0.4cm}
    \noindent \textbf{Keywords:} xxxxxxxxxxxxxxxx

%---------------------------------------------------------
\end{resumoumacoluna}

%%%%%%%%%%%%%%%%%%%%%%%%%%%%%%%%%%%%%%


\end{document}

Página 2 de 2